\section{Дискретные случайные величины}

\subsection{Определение и закон распределения}

\begin{definition}
\textbf{Случайной величиной} $X$ называется числовая функция, ставящая в соответствие каждому элементарному исходу $\omega \in \Omega$ некоторое действительное число $X(\omega)$.
\end{definition}

Случайная величина называется \textbf{дискретной}, если множество ее возможных значений конечно или счетно. Закон ее распределения задается таблицей (рядом распределения), где перечислены значения $x_i$ и соответствующие им вероятности $p_i = P(X = x_i)$. Сумма всех вероятностей обязана равняться единице: $\sum p_i = 1$.

\subsection{Математическое ожидание и его свойства}

\begin{definition}
\textbf{Математическое ожидание} $E(X)$ дискретной случайной величины — это сумма произведений всех ее возможных значений на их вероятности:
$$E(X) = \sum_{i} x_i p_i.$$
Оно характеризует «центр тяжести» распределения.
\end{definition}

\begin{theorem}[Свойства математического ожидания]
Пусть $X$ и $Y$ — случайные величины, а $c, a, b$ — константы.
\begin{enumerate}
\item $E(c) = c$.
\item $E(aX + b) = a E(X) + b$.
\item $E(XY) = E(X) E(Y)$ (только если $X$ и $Y$ независимы).
\end{enumerate}
\end{theorem}

\begin{proof}
Докажем свойства, опираясь исключительно на определение суммы:

1. Константу $c$ можно рассматривать как дискретную случайную величину, которая принимает единственное значение $c$ с вероятностью $p_1 = 1$. Тогда по определению:
$$E(c) = c \cdot 1 = c.$$

2. Рассмотрим новую случайную величину $Z = aX + b$. Ее значения равны $(a x_i + b)$, и они принимаются с теми же вероятностями $p_i$.
$$E(aX + b) = \sum_{i} (a x_i + b) p_i = \sum_{i} (a x_i p_i + b p_i).$$
Разбиваем сумму на две:
$$\sum_{i} a x_i p_i + \sum_{i} b p_i = a \left( \sum_{i} x_i p_i \right) + b \left( \sum_{i} p_i \right).$$
Так как $\sum x_i p_i = E(X)$, а $\sum p_i = 1$ (условие нормировки), получаем $a E(X) + b \cdot 1 = a E(X) + b$.

3. Доказательство для произведения независимых величин приведено в Главе 6 (Двумерные случайные величины), где формализуется совместное распределение.
\end{proof}

\subsection{Дисперсия и ее свойства}

\begin{definition}
\textbf{Дисперсия} $D(X)$ — это математическое ожидание квадрата отклонения случайной величины от ее математического ожидания:
$$D(X) = E\left((X - E(X))^2\right).$$
\end{definition}

\begin{theorem}[Вычислительная формула дисперсии]
Дисперсия равна разности математического ожидания квадрата величины и квадрата ее математического ожидания: $D(X) = E(X^2) - (E(X))^2$.
\end{theorem}

\begin{proof}
Обозначим для краткости константу $\mu = E(X)$. Тогда:
$$D(X) = E\left((X - \mu)^2\right).$$
Раскроем квадрат под знаком математического ожидания алгебраически: $(X - \mu)^2 = X^2 - 2\mu X + \mu^2$.
Используя свойство линейности математического ожидания, разобьем выражение:
$$D(X) = E(X^2 - 2\mu X + \mu^2) = E(X^2) - E(2\mu X) + E(\mu^2).$$
Поскольку $\mu$ — это число (константа), то $E(2\mu X) = 2\mu E(X)$ и $E(\mu^2) = \mu^2$.
Подставляем обратно $E(X) = \mu$:
$$D(X) = E(X^2) - 2\mu \cdot \mu + \mu^2 = E(X^2) - 2\mu^2 + \mu^2 = E(X^2) - \mu^2 = E(X^2) - (E(X))^2.$$
Теорема доказана.
\end{proof}

\begin{theorem}[Свойства дисперсии]
\begin{enumerate}
\item $D(c) = 0$.
\item $D(X + c) = D(X)$.
\item $D(aX) = a^2 D(X)$.
\end{enumerate}
\end{theorem}

\begin{proof}
\begin{enumerate}
\item $D(c) = E(c^2) - (E(c))^2 = c^2 - c^2 = 0$. 
\item $D(X + c) = E(((X+c) - E(X+c))^2) = E((X + c - E(X) - c)^2) = E((X - E(X))^2) = D(X)$. Сдвиг распределения не меняет его разброс. 
\item $D(aX) = E((aX - E(aX))^2) = E((aX - aE(X))^2) = E(a^2 (X - E(X))^2)$. Выносим константу $a^2$ за знак $E$: получаем $a^2 E((X - E(X))^2) = a^2 D(X)$.
\end{enumerate}
\end{proof}